% !PRINTED 170
\weberpage{170}
Let $F_1, F_2, \ldots, F_\mu$ be a basis of the system $\mathfrak{J}$ determined according to II., and let $F$ be any other non-constant invariant of $S$. Then the forms $A_1, A_2, \ldots, A_\mu$ can be so determined that
\weq{2}{F = A_1F_1 + A_2F_2 + \cdots + A_\mu F_\mu}
holds. If one applies to this identical equation all the substitutions of the group $S$, then by hypothesis $F, F_1, F_2, \ldots F_\mu$ remain unchanged, while $A_\varkappa$ may pass over into $A_\varkappa, A_\varkappa', A_\varkappa'' \ldots$ If one forms the sum of the equations thus derived from (2) and sets, where $m$ denotes the degree of the group $S$,
\weq{3}{m\Phi_\varkappa = A_\varkappa + A_\varkappa' + A_\varkappa'' + \cdots,}
then it follows that:
\weq{4}{F = \Phi_1F_1 + \Phi_2F_2 + \cdots + \Phi_\mu F_\mu.}
But the functions $\Phi_1, \Phi_2, \ldots \Phi_\mu$ are by \S.\,40 invariants of $S$, and with this, according to III., the finiteness of the system $\mathfrak{J}$ is proved.

The same conclusion can also be applied to the relative invariants, as follows:

We denote by $F(x)$ the whole system of the functions which satisfy the conditions \S.\,40, (3):
\[
F[A(x)] = \alpha F(x), \quad F[B(x)] = \beta F(x) \ldots
\]
with fixed factors $\alpha, \beta, \ldots$, which, as we have seen, are roots of unity.

By \gesperrt{Hilbert's} theorem a special system of such functions $F_1, F_2, \ldots, F_m$ can be selected in such a way that one can set
\weq{5}{F = A_1F_1 + A_2F_2 + \cdots + A_mF_m}
where $A_1, A_2, \ldots, A_m$ are forms of the variables $(x)$. If one treats this formula just as the formula (2), by applying to it the substitutions of the group $S$ and then forming the sum, one obtains, corresponding to the formula (4):
\[
F = \Phi_1F_1 + \Phi_2F_2 + \cdots + \Phi_mF_m,
\]
in which the coefficients $\Phi_1, \Phi_2, \ldots, \Phi_m$ are \gesperrt{absolute invariants}.

By way of completion it must still be added that inhomogeneous functions of the variables can be invariants only when their individual homogeneous parts are invariants.
